% GENERATED FILE. Edit canonical lessons, then run npm run generate:books.
% canonical-source-hash: fnv1a64:3e5ea75fbe423097
% canonical-lessons: SA-C319-vijayah, SA-C319-parajayah, SA-C319-bandi, SA-C319-apasarati, SA-C319-abhibhavati

\chapter{Victory and Defeat}
\label{ch:sa-victory-and-defeat}
\hlchaptermodality{319}

\begin{chapteropening}
\textbf{By the end of this chapter:} \emph{I can say who won a battle and what became of the losers.}

Complete the last lesson of chapter 319: I can say who won a battle and what became of the losers.
\end{chapteropening}

\section[\texorpdfstring{\sk{विजयः} (vijayaḥ)}{विजयः (vijayaḥ)}]{\sk{विजयः} (vijayaḥ) — a victory}

\label{lesson:SA-C319-vijayah}



\begin{quote}
Before the new one: say the Sanskrit for sets out, then the Sanskrit for fights.
\end{quote}

\subsection*{You'll want to know: \sk{विजयः}}
\textbf{\sk{विजयः}} — \emph{vijayaḥ} — \textquotedblleft{}a victory\textquotedblright{}.

\begin{quote}
\sk{अन्ते} \sk{सेनायाः} \sk{विजयः} \sk{भवति।}
\end{quote}

\begin{quote}
\emph{ante senāyāḥ vijayaḥ bhavati.}
\end{quote}

\begin{quote}
In the end the army wins the victory.
\end{quote}

A masculine noun.

\begin{cousinweb}
\textbf{\sk{विजयः}} (\emph{vijayaḥ}) is built from \emph{vi} and the root √ji, \textquotedblleft{}to conquer, win\textquotedblright{}, as in \emph{jayati}; MW gives \emph{vijaya} as a contest for victory, victory, conquest (MW).
\end{cousinweb}

\begin{grammarlens}[title={What you've built}]
One more word for telling what happened.
\end{grammarlens}

\subsection*{Your turn}
\begin{itemize}
\raggedright
  \item \emph{Say it:} \emph{vijayaḥ}
  \item \emph{Say it:} \emph{vijayaḥ}, once more
  \item \emph{Say it:} \emph{ante senāyāḥ vijayaḥ bhavati.} — In the end the army wins the victory.
  \item \emph{Recall:} say the Sanskrit for sets out, then the Sanskrit for fights, then say \emph{vijayaḥ} again
\end{itemize}

\begin{tcolorbox}[breakable,colback=teal!4,colframe=teal!35!black,title={Before you move on}]
What does \sk{विजयः} mean? (\textquotedblleft{}A victory\textquotedblright{}.) Say it once more.
\end{tcolorbox}
\section[\texorpdfstring{\sk{पराजयः} (parājayaḥ)}{पराजयः (parājayaḥ)}]{\sk{पराजयः} (parājayaḥ) — a defeat}

\label{lesson:SA-C319-parajayah}



\begin{quote}
Before the new one: say the Sanskrit for fights, then the Sanskrit for a victory.
\end{quote}

\subsection*{You'll want to know: \sk{पराजयः}}
\textbf{\sk{पराजयः}} — \emph{parājayaḥ} — \textquotedblleft{}a defeat\textquotedblright{}.

\begin{quote}
\sk{शत्रोः} \sk{पराजयः} \sk{स्पष्टः} \sk{अस्ति।}
\end{quote}

\begin{quote}
\emph{śatroḥ parājayaḥ spaṣṭaḥ asti.}
\end{quote}

\begin{quote}
The enemy's defeat is clear.
\end{quote}

A masculine noun.

\begin{cousinweb}
\textbf{\sk{पराजयः}} (\emph{parājayaḥ}) is built from \emph{parā}, \textquotedblleft{}away\textquotedblright{}, and \emph{jaya}, \textquotedblleft{}victory\textquotedblright{}, from the root √ji; MW gives \emph{parājaya} as the being deprived of something, loss, defeat (MW).
\end{cousinweb}

\begin{grammarlens}[title={What you've built}]
One more word for telling what happened.
\end{grammarlens}

\subsection*{Your turn}
\begin{itemize}
\raggedright
  \item \emph{Say it:} \emph{parājayaḥ}
  \item \emph{Say it:} \emph{parājayaḥ}, once more
  \item \emph{Say it:} \emph{śatroḥ parājayaḥ spaṣṭaḥ asti.} — The enemy's defeat is clear.
  \item \emph{Recall:} say the Sanskrit for fights, then the Sanskrit for a victory, then say \emph{parājayaḥ} again
\end{itemize}

\begin{tcolorbox}[breakable,colback=teal!4,colframe=teal!35!black,title={Before you move on}]
What does \sk{पराजयः} mean? (\textquotedblleft{}A defeat\textquotedblright{}.) Say it once more.
\end{tcolorbox}
\section[\texorpdfstring{\sk{बन्दी} (bandī)}{बन्दी (bandī)}]{\sk{बन्दी} (bandī) — a captive; a prisoner of war}

\label{lesson:SA-C319-bandi}



\begin{quote}
Before the new one: say the Sanskrit for a victory, then the Sanskrit for a defeat.
\end{quote}

\subsection*{You'll want to know: \sk{बन्दी}}
\textbf{\sk{बन्दी}} — \emph{bandī} — \textquotedblleft{}a captive; a prisoner of war\textquotedblright{}.

\begin{quote}
\sk{एकः} \sk{बन्दी} \sk{कारागारे} \sk{अस्ति।}
\end{quote}

\begin{quote}
\emph{ekaḥ bandī kārāgāre asti.}
\end{quote}

\begin{quote}
One captive is in prison.
\end{quote}

A masculine noun.

\begin{cousinweb}
\textbf{\sk{बन्दी}} (\emph{bandī}) MW gives this \emph{bandin} (also written \emph{vandin}) as a prisoner, captive, slave, without a derivation in the entry, so its origin is left uncertain here (MW).
\end{cousinweb}

\begin{grammarlens}[title={What you've built}]
One more word for telling what happened.
\end{grammarlens}

\subsection*{Your turn}
\begin{itemize}
\raggedright
  \item \emph{Say it:} \emph{bandī}
  \item \emph{Say it:} \emph{bandī}, once more
  \item \emph{Say it:} \emph{ekaḥ bandī kārāgāre asti.} — One captive is in prison.
  \item \emph{Recall:} say the Sanskrit for a victory, then the Sanskrit for a defeat, then say \emph{bandī} again
\end{itemize}

\begin{tcolorbox}[breakable,colback=teal!4,colframe=teal!35!black,title={Before you move on}]
What does \sk{बन्दी} mean? (\textquotedblleft{}A captive; a prisoner of war\textquotedblright{}.) Say it once more.
\end{tcolorbox}
\section[\texorpdfstring{\sk{अपसरति} (apasarati)}{अपसरति (apasarati)}]{\sk{अपसरति} (apasarati) — retreats; moves away}

\label{lesson:SA-C319-apasarati}



\begin{quote}
Before the new one: say the Sanskrit for a defeat, then the Sanskrit for a captive.
\end{quote}

\subsection*{You'll want to know: \sk{अपसरति}}
\textbf{\sk{अपसरति}} — \emph{apasarati} — \textquotedblleft{}retreats; moves away\textquotedblright{}.

\begin{quote}
\sk{शत्रोः} \sk{सेना} \sk{शीघ्रम्} \sk{अपसरति।}
\end{quote}

\begin{quote}
\emph{śatroḥ senā śīghram apasarati.}
\end{quote}

\begin{quote}
The enemy's army quickly retreats.
\end{quote}

Past (it …): \sk{सेना} \sk{शीघ्रम्} \sk{अपसृतवती}, \emph{senā śīghram apasṛtavatī}.

\begin{cousinweb}
\textbf{\sk{अपसरति}} (\emph{apasarati}) is built from \emph{apa}, \textquotedblleft{}away\textquotedblright{}, and the root √sṛ, \textquotedblleft{}to go, run\textquotedblright{}; MW gives \emph{apa-sṛ} as to slip off from, go away, retreat (MW).
\end{cousinweb}

\begin{grammarlens}[title={What you've built}]
One more word for telling what happened.
\end{grammarlens}

\subsection*{Your turn}
\begin{itemize}
\raggedright
  \item \emph{Say it:} \emph{apasarati}
  \item \emph{Say it:} \emph{apasarati}, once more
  \item \emph{Say it:} \emph{śatroḥ senā śīghram apasarati.} — The enemy's army quickly retreats.
  \item \emph{Recall:} say the Sanskrit for a defeat, then the Sanskrit for a captive, then say \emph{apasarati} again
\end{itemize}

\begin{tcolorbox}[breakable,colback=teal!4,colframe=teal!35!black,title={Before you move on}]
What does \sk{अपसरति} mean? (\textquotedblleft{}Retreats; moves away\textquotedblright{}.) Say it once more.
\end{tcolorbox}
\section[\texorpdfstring{\sk{अभिभवति} (abhibhavati)}{अभिभवति (abhibhavati)}]{\sk{अभिभवति} (abhibhavati) — overpowers}

\label{lesson:SA-C319-abhibhavati}



\begin{quote}
Before the new one: say the Sanskrit for a captive, then the Sanskrit for retreats.
\end{quote}

\subsection*{You'll want to know: \sk{अभिभवति}}
\textbf{\sk{अभिभवति}} — \emph{abhibhavati} — \textquotedblleft{}overpowers\textquotedblright{}.

\begin{quote}
\sk{वीरः} \sk{राजा} \sk{शत्रुम्} \sk{अभिभवति।}
\end{quote}

\begin{quote}
\emph{vīraḥ rājā śatrum abhibhavati.}
\end{quote}

\begin{quote}
The brave king overpowers the enemy.
\end{quote}

Past (he …): \sk{सः} \sk{शत्रुम्} \sk{अभिभूतवान्}, \emph{saḥ śatrum abhibhūtavān}.

\begin{cousinweb}
\textbf{\sk{अभिभवति}} (\emph{abhibhavati}) is built from \emph{abhi}, \textquotedblleft{}over, against\textquotedblright{}, and the root √bhū, \textquotedblleft{}to be, become\textquotedblright{}, as in \emph{bhavati}; MW gives \emph{abhi-bhū} as to overcome, overpower, defeat (MW).
\end{cousinweb}

\begin{grammarlens}[title={What you've built}]
One more word for telling what happened.
\end{grammarlens}

\subsection*{Your turn}
\begin{itemize}
\raggedright
  \item \emph{Say it:} \emph{abhibhavati}
  \item \emph{Say it:} \emph{abhibhavati}, once more
  \item \emph{Say it:} \emph{vīraḥ rājā śatrum abhibhavati.} — The brave king overpowers the enemy.
  \item \emph{Recall:} say the Sanskrit for a captive, then the Sanskrit for retreats, then say \emph{abhibhavati} again
\end{itemize}

\begin{tcolorbox}[breakable,colback=teal!4,colframe=teal!35!black,title={Before you move on}]
What does \sk{अभिभवति} mean? (\textquotedblleft{}Overpowers\textquotedblright{}.) Say it once more.
\end{tcolorbox}
